% =====================================================================
\chapter{Reducción al absurdo y demostración por casos}
\label{tema:4}
% =====================================================================

Algunas afirmaciones no tienen la forma de un condicional, o la tienen pero ni la hipótesis ni la negación de la tesis dan un punto de partida útil. Para ellas se dispone de la reducción al absurdo: suponer lo contrario de lo que se quiere probar y llegar a una contradicción. Este tema la presenta con dos de los teoremas más antiguos de las matemáticas, la irracionalidad de $\sqrt2$ y la infinitud de los números primos, y estudia también la demostración por casos. Se siguen el capítulo~6 y las §4.4–4.5 de~\cite{hammack}, las §1.7–1.8 de~\cite{mcs}, la §3.5 de~\cite{velleman} y el primer \emph{problem set} de CS103~\cite{cs103}.

\section{Reducción al absurdo}

\begin{definicion}[Reducción al absurdo]\label{def:absurdo}
Para demostrar una proposición $P$ por \emph{reducción al absurdo} (o por contradicción) se supone $\neg P$ y se deduce una contradicción, es decir, una proposición de la forma $C\wedge\neg C$. Se concluye $P$.
\end{definicion}

La justificación es que $(\neg P\Rightarrow(C\wedge\neg C))\Rightarrow P$ es una tautología: si $P$ fuera falsa, $\neg P$ sería verdadera, y el condicional $\neg P\Rightarrow(C\wedge\neg C)$, con hipótesis verdadera y tesis falsa, sería falso. Para un enunciado condicional $P\Rightarrow Q$, la negación es $P\wedge\neg Q$ (\cref{prop:leyes}(vi)): se suponen a la vez la hipótesis y la negación de la tesis~\cite[§6.2]{hammack}.

\begin{ejemplo}\label{ej:no-mayor}
No existe un número entero mayor que todos los demás. Supóngase, por reducción al absurdo, que existe $M\in\Z$ con $n\le M$ para todo $n\in\Z$. Tomando $n=M+1\in\Z$ se obtiene $M+1\le M$, es decir, $1\le 0$, que es falso. Por tanto, tal $M$ no existe.
\end{ejemplo}

\begin{ejemplo}\label{ej:no-solucion}
No existen enteros $a$ y $b$ con $18a+6b=1$. Supóngase que sí existen. Entonces $1=18a+6b=2(9a+3b)$ sería par, mientras que $1=2\cdot 0+1$ es impar; ningún entero es par e impar a la vez (\cref{teo:division}). Contradicción.
\end{ejemplo}

\section{La irracionalidad de \texorpdfstring{$\sqrt2$}{raíz de 2}}

Un número real es \emph{racional} si se escribe como $a/b$ con $a,b\in\Z$ y $b\neq 0$, e \emph{irracional} en caso contrario. Todo racional puede escribirse en \emph{forma irreducible}, $a/b$ con $b>0$ y sin divisores comunes mayores que~1: basta dividir $a$ y $b$ por su máximo común divisor. (Alternativamente, se toma la representación con $b$ positivo mínimo, que existe por el principio del buen orden del \cref{tema:5}.)

\begin{teorema}\label{teo:raiz2}
$\sqrt2$ es irracional.
\end{teorema}

\begin{proof}
Supóngase, por reducción al absurdo, que $\sqrt2$ es racional, y escríbase $\sqrt2=a/b$ con $a,b\in\Z$, $b>0$ y la fracción irreducible. Elevando al cuadrado, $2=a^{2}/b^{2}$, es decir,
\[
a^{2}=2b^{2}.
\]
Así, $a^{2}$ es par, y por la \cref{prop:cuadrado-par} $a$ es par: $a=2c$ con $c\in\Z$. Sustituyendo, $4c^{2}=2b^{2}$, luego $b^{2}=2c^{2}$ es par, y de nuevo $b$ es par. Entonces $2$ divide a $a$ y a $b$, en contra de que $a/b$ sea irreducible. Por tanto, $\sqrt2$ no es racional.
\end{proof}

La demostración sigue la del teorema~1.8.1 de~\cite{mcs}. Su estructura se repite en los \cref{eje:4-raiz3,eje:4-raiz-cubica,eje:4-log}: suponer una fracción irreducible, obtener una ecuación entre enteros y encontrar un divisor común.

\section{La infinitud de los números primos}

\begin{lema}\label{lem:divisor-primo}
Todo entero $n\ge 2$ tiene al menos un divisor primo.
\end{lema}

\begin{proof}
El conjunto $D=\{d\in\Z: 2\le d\le n,\ d\divides n\}$ es finito y no vacío, pues $n\in D$. Sea $p$ su menor elemento. Si $p$ no fuera primo, tendría un divisor $e$ con $2\le e<p$; por la \cref{prop:divisibilidad}(i), $e\divides n$, luego $e\in D$ y $e<p$, en contra de la minimalidad de $p$. Por tanto, $p$ es un divisor primo de $n$.
\end{proof}

\begin{teorema}[Euclides]\label{teo:euclides}
Existen infinitos números primos.
\end{teorema}

\begin{proof}
Supóngase, por reducción al absurdo, que hay solo un número finito de primos, $p_1,p_2,\dots,p_k$. Sea
\[
N=p_1p_2\cdots p_k+1.
\]
Como $N\ge 2+1\ge 2$, el \cref{lem:divisor-primo} da un primo $p$ que divide a $N$, y por la hipótesis $p=p_i$ para algún $i$. Entonces $p_i$ divide a $N$ y a $p_1\cdots p_k$, y por la \cref{prop:divisibilidad}(ii) divide a su diferencia, $N-p_1\cdots p_k=1$. Pero ningún primo divide a~1, porque $p_i\ge 2$. Contradicción.
\end{proof}

\begin{observacion}
La prueba \emph{no} afirma que $N$ sea primo. Con los seis primeros primos, $N=2\cdot 3\cdot 5\cdot 7\cdot 11\cdot 13+1=30\,031=59\cdot 509$ es compuesto; lo que la prueba garantiza es que sus factores primos no están en la lista (\cref{eje:4-euclides-n}).
\end{observacion}

\section{Demostración por casos}

\begin{definicion}[Demostración por casos]\label{def:casos}
Para demostrar $Q$ a partir de $P_1\vee P_2\vee\dots\vee P_k$ se demuestran por separado $P_1\Rightarrow Q$, \dots, $P_k\Rightarrow Q$. Es válido porque $(P_1\vee\dots\vee P_k)\Rightarrow Q\eqlog(P_1\Rightarrow Q)\wedge\dots\wedge(P_k\Rightarrow Q)$, generalización del \cref{eje:1-equivalencias}(c).
\end{definicion}

Lo esencial es que los casos sean \emph{exhaustivos}: su disyunción debe ser verdadera. Las divisiones más habituales son «$n$ par o impar», los restos módulo $m$ y el signo de un número real~\cite[§4.4]{hammack}.

\begin{proposicion}\label{prop:consecutivos}
Para todo $n\in\Z$, el número $n^{2}+n$ es par.
\end{proposicion}

\begin{proof}
Todo entero es par o impar. \emph{Caso 1:} $n=2k$. Entonces $n^{2}+n=4k^{2}+2k=2(2k^{2}+k)$ es par. \emph{Caso 2:} $n=2k+1$. Entonces $n^{2}+n=(2k+1)(2k+2)=2(2k+1)(k+1)$ es par. En ambos casos $n^{2}+n$ es par.
\end{proof}

\begin{proposicion}\label{prop:cuadrado-mod4}
Para todo $n\in\Z$, $n^{2}\equiv 0\pmod 4$ o $n^{2}\equiv 1\pmod 4$.
\end{proposicion}

\begin{proof}
Si $n=2k$, $n^{2}=4k^{2}\equiv 0$. Si $n=2k+1$, $n^{2}=4(k^{2}+k)+1\equiv 1\pmod 4$.
\end{proof}

\begin{teorema}[Desigualdad triangular]\label{teo:triangular}
Para todos $x,y\in\R$, $\abs{x+y}\le\abs{x}+\abs{y}$.
\end{teorema}

\begin{proof}
Recuérdese que $\abs{t}=t$ si $t\ge 0$ y $\abs{t}=-t$ si $t<0$; en ambos casos $-\abs{t}\le t\le\abs{t}$. Se distinguen dos casos según el signo de $x+y$.
\emph{Caso 1:} $x+y\ge 0$. Entonces $\abs{x+y}=x+y\le\abs{x}+\abs{y}$.
\emph{Caso 2:} $x+y<0$. Entonces $\abs{x+y}=-x-y=(-x)+(-y)\le\abs{x}+\abs{y}$.
\end{proof}

Cuando dos casos se tratan con el mismo argumento intercambiando los papeles de las variables, se escribe «sin pérdida de generalidad, supóngase $x\le y$» y se trata solo uno~\cite[§4.5]{hammack}. Debe usarse solo cuando la simetría es real.

\begin{observacion}[Errores frecuentes]\label{obs:4-errores}
(a) El «sándwich de contradicción»: suponer $\neg P$, demostrar $P$ directamente sin usar la suposición y declarar una contradicción. Esa prueba es directa y la reducción al absurdo sobra~\cite{cs103}.
(b) Negar mal: la negación de «$\forall n$, $n$ es impar» no es «$\forall n$, $n$ es par» (\cref{eje:2-critica}).
(c) Casos no exhaustivos: dividir un real en «positivo» y «negativo» olvida el cero.
(d) En una reducción al absurdo, señalar con claridad qué se supone, cuál es la contradicción y qué se concluye.
\end{observacion}

\begin{observacion}[Conexión con el curso]
El módulo~4 demuestra por reducción al absurdo, con el teorema de Taylor, que en un mínimo local la hessiana es semidefinida positiva, y el módulo~6 usa un argumento por contradicción sobre los códigos prefijo. Las demostraciones por casos aparecen en cuanto intervienen funciones definidas a trozos: el valor absoluto en las cotas de error del módulo~17, la función ReLU, $\max(0,x)$, y su derivada en los módulos~9 y~11, y la desigualdad triangular en las normas del módulo~2.
\end{observacion}

\begin{lecturas}
\item Hammack~\cite{hammack}, cap.~6 «Proof by Contradiction» (§6.1–6.4) y §4.4 «Using Cases», §4.5 «Treating Similar Cases».
\item Lehman, Leighton y Meyer~\cite{mcs}, §1.7 «Proof by Cases» y §1.8 «Proof by Contradiction» (teorema~1.8.1).
\item Velleman~\cite{velleman}, §3.5 «Proofs Involving Disjunctions».
\item Stanford CS103~\cite{cs103}, clase~2 «Indirect Proofs».
\end{lecturas}

\section*{Ejercicios del tema 4}
\addcontentsline{toc}{section}{Ejercicios del tema 4}

\begin{ejercicio}{1}{eje:4-suponer}
Escriba qué se supone al comenzar una reducción al absurdo de: (a) «$\sqrt5$ es irracional»; (b) «para todo $n\in\N$, $n^{2}+1$ no es múltiplo de~3»; (c) «si $a$ y $b$ son enteros con $ab$ impar, entonces $a$ y $b$ son impares»; (d) «existe un primo mayor que $10^{100}$».
\end{ejercicio}

\begin{ejercicio}{3}{eje:4-raiz3}
Demuestre que, si $n\in\Z$ y $3\divides n^{2}$, entonces $3\divides n$ (por contrarrecíproco y casos). Deduzca que $\sqrt3$ es irracional.
\fuente{adaptado de Hammack~\cite{hammack}, cap.~6 «Proof by Contradiction», y de Lehman, Leighton y Meyer~\cite{mcs}, teorema~1.8.1.}
\end{ejercicio}

\begin{ejercicio}{3}{eje:4-raiz-cubica}
Demuestre que $\sqrt[3]{2}$ es irracional.
\end{ejercicio}

\begin{ejercicio}{3}{eje:4-log}
Demuestre que $\log_{2}3$ es irracional.
\fuente{adaptado de Lehman, Leighton y Meyer~\cite{mcs}, §1.8 «Proof by Contradiction» (problemas de la sección).}
\end{ejercicio}

\begin{ejercicio}{3}{eje:4-racional-irracional}
Sean $a$ un número racional y $b$ un número irracional. Demuestre que $a+b$ es irracional y que, si $a\neq 0$, también $ab$ es irracional. ¿Qué ocurre con $ab$ si $a=0$?
\end{ejercicio}

\begin{ejercicio}{3}{eje:4-diofantica}
Demuestre que no existen enteros $a$ y $b$ con $21a+30b=1$. ¿Existen enteros con $21a+30b=3$?
\fuente{adaptado de Hammack~\cite{hammack}, cap.~6 «Proof by Contradiction».}
\end{ejercicio}

\begin{ejercicio}{3}{eje:4-mod4}
Demuestre que, para todo $n\in\Z$, $n^{2}+2$ no es múltiplo de~4. Demuestre también que ningún entero de la forma $4k+3$ es suma de dos cuadrados de enteros.
\end{ejercicio}

\begin{ejercicio}{3}{eje:4-24}
Sea $p$ un número primo mayor que~3. Demuestre que $24\divides p^{2}-1$.
\end{ejercicio}

\begin{ejercicio}{2}{eje:4-maximo}
Demuestre por casos que, para todos $x,y\in\R$, $\max(x,y)=\dfrac{x+y+\abs{x-y}}{2}$ y $\min(x,y)=\dfrac{x+y-\abs{x-y}}{2}$. Compruebe las fórmulas con $(x,y)=(3,-5)$.
\end{ejercicio}

\begin{ejercicio}{2}{eje:4-restos-cuadrados}
Construya, por casos según el resto de $n$, las tablas de $n^{2}$ módulo~3, módulo~5 y módulo~8. Deduzca: (a) $n^{2}+1$ nunca es múltiplo de~3; (b) ningún cuadrado es congruente con $2$ ni con $3$ módulo~5; (c) el cuadrado de todo entero impar es congruente con~$1$ módulo~8.
\end{ejercicio}

\begin{ejercicio}{3}{eje:4-triangular-inversa}
Demuestre que, para todos $x,y\in\R$, $\bigl\lvert\abs{x}-\abs{y}\bigr\rvert\le\abs{x-y}$. (Indicación: aplique el \cref{teo:triangular} a $x=(x-y)+y$ y a $y=(y-x)+x$.)
\end{ejercicio}

\begin{ejercicio}{3}{eje:4-producto-par}
Sean $m,n\in\Z$. Demuestre por reducción al absurdo que, si $mn$ es par y $m$ es impar, entonces $n$ es par.
\fuente{adaptado de Stanford CS103~\cite{cs103}, \emph{Problem Set}~1, problema~6 «Writing Proofs by Contradiction».}
\end{ejercicio}

\begin{ejercicio}{3}{eje:4-cubos}
(a) Demuestre que todo múltiplo de~6 es suma de cuatro cubos de enteros. (b) Demuestre que $x^{3}\equiv x\pmod 6$ para todo $x\in\Z$. (c) Deduzca que todo entero es suma de cinco cubos de enteros y escriba $7$ de esa forma.
\fuente{adaptado de Stanford CS103~\cite{cs103}, \emph{Problem Set}~1, problema~10 «Sums of Cubes».}
\end{ejercicio}

\begin{ejercicio}{2}{eje:4-euclides-n}
Para $k=1,\dots,6$, sea $N_k=p_1\cdots p_k+1$, donde $p_1,\dots,p_k$ son los $k$ primeros primos. Calcule $N_k$, decida cuáles son primos y factorice $N_6$. Explique por qué el resultado no contradice la demostración del \cref{teo:euclides}.
\end{ejercicio}

\begin{ejercicio}{3}{eje:4-4k3}
Demuestre que existen infinitos primos de la forma $4k+3$. (Indicación: suponga que $p_1,\dots,p_m$ son todos, considere $N=4p_1\cdots p_m-1$ y demuestre que algún divisor primo de $N$ es de la forma $4k+3$.)
\end{ejercicio}

\begin{ejercicio}{3}{eje:4-raiz-primo}
Sea $p$ un número primo. Admitiendo el \emph{lema de Euclides} (si $p\divides ab$, entonces $p\divides a$ o $p\divides b$), demuestre que $\sqrt{p}$ es irracional. ¿En qué paso falla la demostración si se intenta con $\sqrt4$?
\end{ejercicio}
